\documentclass[10pt,a4paper]{amsart}
\usepackage{amsmath,amssymb,amsthm,graphicx}
% Retained minimal preamble for the Stacks Project chapter more-morphisms.tex (More on Morphisms).
% Verbatim excerpt of the chapter preamble: Xy-pic setup (native preamble lines 14-18) and the native
% theorem declarations with the mathematical macros (native preamble lines 154-217). Copied mechanically;
% nothing added, removed or redefined. The wrapper class and the amsmath/amssymb/amsthm/graphicx base
% packages are supplied by the assembler, so the class/externaldocument part of the native preamble is omitted.
% For commutative diagrams we use Xy-pic
\usepackage[all]{xy}

% We use 2cell for 2-commutative diagrams.
\xyoption{2cell}
% Theorem environments.
%
\theoremstyle{plain}
\newtheorem{theorem}[subsection]{Theorem}
\newtheorem{proposition}[subsection]{Proposition}
\newtheorem{lemma}[subsection]{Lemma}

\theoremstyle{definition}
\newtheorem{definition}[subsection]{Definition}
\newtheorem{example}[subsection]{Example}
\newtheorem{exercise}[subsection]{Exercise}
\newtheorem{situation}[subsection]{Situation}

\theoremstyle{remark}
\newtheorem{remark}[subsection]{Remark}
\newtheorem{remarks}[subsection]{Remarks}

\numberwithin{equation}{subsection}

% Macros
%
\def\lim{\mathop{\mathrm{lim}}\nolimits}
\def\colim{\mathop{\mathrm{colim}}\nolimits}
\def\Spec{\mathop{\mathrm{Spec}}}
\def\Hom{\mathop{\mathrm{Hom}}\nolimits}
\def\Ext{\mathop{\mathrm{Ext}}\nolimits}
\def\SheafHom{\mathop{\mathcal{H}\!\mathit{om}}\nolimits}
\def\SheafExt{\mathop{\mathcal{E}\!\mathit{xt}}\nolimits}
\def\Sch{\mathit{Sch}}
\def\Mor{\mathop{\mathrm{Mor}}\nolimits}
\def\Ob{\mathop{\mathrm{Ob}}\nolimits}
\def\Sh{\mathop{\mathit{Sh}}\nolimits}
\def\NL{\mathop{N\!L}\nolimits}
\def\CH{\mathop{\mathrm{CH}}\nolimits}
\def\proetale{{pro\text{-}\acute{e}tale}}
\def\etale{{\acute{e}tale}}
\def\QCoh{\mathit{QCoh}}
\def\Ker{\mathop{\mathrm{Ker}}}
\def\Im{\mathop{\mathrm{Im}}}
\def\Coker{\mathop{\mathrm{Coker}}}
\def\Coim{\mathop{\mathrm{Coim}}}

% Boxtimes
%
\DeclareMathSymbol{\boxtimes}{\mathbin}{AMSa}{"02}

%
% Macros for moduli stacks/spaces
%
\def\QCohstack{\mathcal{QC}\!\mathit{oh}}
\def\Cohstack{\mathcal{C}\!\mathit{oh}}
\def\Spacesstack{\mathcal{S}\!\mathit{paces}}
\def\Quotfunctor{\mathrm{Quot}}
\def\Hilbfunctor{\mathrm{Hilb}}
\def\Curvesstack{\mathcal{C}\!\mathit{urves}}
\def\Polarizedstack{\mathcal{P}\!\mathit{olarized}}
\def\Complexesstack{\mathcal{C}\!\mathit{omplexes}}
% \Pic is the operator that assigns to X its picard group, usage \Pic(X)
% \Picardstack_{X/B} denotes the Picard stack of X over B
% \Picardfunctor_{X/B} denotes the Picard functor of X over B
\def\Pic{\mathop{\mathrm{Pic}}\nolimits}
\def\Picardstack{\mathcal{P}\!\mathit{ic}}
\def\Picardfunctor{\mathrm{Pic}}
\def\Deformationcategory{\mathcal{D}\!\mathit{ef}}
\usepackage{cleveref}
\crefname{originallemma}{original lemma}{original lemmas}
\Crefname{originallemma}{Original lemma}{Original lemmas}
\makeatletter
\newlabel{lemma-affineness-of-large-open}{{0F3R}{source}{}{sourcepointer-1}{}}
\newlabel{lemma-affineness-of-large-open@cref}{{[originallemma][0][]0F3R}{[1][0][]source}{}{}{}}
\makeatother
\crefname{originallemma}{original lemma}{original lemmas}
\Crefname{originallemma}{Original lemma}{Original lemmas}
\makeatletter
\newlabel{lemma-max-int-finite}{{0F3L}{source}{}{sourcepointer-2}{}}
\newlabel{lemma-max-int-finite@cref}{{[originallemma][0][]0F3L}{[1][0][]source}{}{}{}}
\makeatother
\crefname{originallemma}{original lemma}{original lemmas}
\Crefname{originallemma}{Original lemma}{Original lemmas}
\makeatletter
\newlabel{lemma-max-int-w}{{0F3K}{source}{}{sourcepointer-3}{}}
\newlabel{lemma-max-int-w@cref}{{[originallemma][0][]0F3K}{[1][0][]source}{}{}{}}
\makeatother
\crefname{originallemma}{original lemma}{original lemmas}
\Crefname{originallemma}{Original lemma}{Original lemmas}
\makeatletter
\newlabel{lemma-weighting-base-change}{{0F3B}{source}{}{sourcepointer-4}{}}
\newlabel{lemma-weighting-base-change@cref}{{[originallemma][0][]0F3B}{[1][0][]source}{}{}{}}
\makeatother
\crefname{originalreference}{original reference}{original references}
\Crefname{originalreference}{Original reference}{Original references}
\makeatletter
\newlabel{limits-lemma-affine}{{01ZT}{source}{}{sourcepointer-5}{}}
\newlabel{limits-lemma-affine@cref}{{[originalreference][0][]01ZT}{[1][0][]source}{}{}{}}
\makeatother
\crefname{originalreference}{original reference}{original references}
\Crefname{originalreference}{Original reference}{Original references}
\makeatletter
\newlabel{morphisms-lemma-affine-permanence}{{01SG}{source}{}{sourcepointer-6}{}}
\newlabel{morphisms-lemma-affine-permanence@cref}{{[originalreference][0][]01SG}{[1][0][]source}{}{}{}}
\makeatother
\title{The affine stratification number under a positive weighting of a surjective quasi-finite morphism}
\author{The Stacks Project Authors}
\date{stacks-project:more-morphisms.tex:local-snapshot-20261011; GNU FDL 1.2 or later}
\begin{document}
\maketitle
\noindent\textit{Editorial note.} Counted span is the author's own proof environment, source lines 23089-23118, immediately following the statement at 23079-23087.\par\medskip
\noindent\textit{Editorial note.} The footnote on 'separated' (source line 23082: it suffices that the diagonal of Y is affine) is retained verbatim inside the counted statement span.\par\medskip
\noindent\textit{Editorial note.} The same-chapter pointer family (the max-int lemmas and the weighting base-change lemma) consists of line-indexed stubs into the pinned source.\par\medskip
\noindent\textit{Allowed standard prerequisites (editorial).}\begin{itemize}
\item Separated target makes quasi-finite morphisms affine. Morphisms, Lemma morphisms-lemma-affine-permanence, cited at proof lines 23092-23093. Gives that X -> Y is affine once Y is separated, the standing input for every 'affine over X is affine' step.
\item Maximum of the weighted integral and finiteness over its stratum. Lemmas lemma-max-int-w (line 23095) and lemma-max-int-finite (lines 23098-23099) of this chapter. Supplies the induction parameter d and that f\textasciicircum{}\{-1\}(Y\_d) -> Y\_d is finite.
\item Affineness of the maximum stratum. Lemma lemma-affineness-of-large-open of this chapter, cited at proof lines 23099-23102. Y\_d meets every affine open W of Y affinely, hence Y\_d -> Y is an affine morphism of schemes.
\item Affineness from an affine surjective cover. Limits, Lemma limits-lemma-affine, cited at proof lines 23108-23109. Y\_d is affine because it is surjected upon by the affine scheme f\textasciicircum{}\{-1\}(Y\_d).
\item Weightings are stable under base change. Lemma lemma-weighting-base-change of this chapter, cited at proof lines 23119-23121. Lets the induction start again on the reduced closed part Y' with the restricted weighting w'.
\end{itemize}\medskip
\ifdefined\hypertarget\hypertarget{sourcepointer-1}{}\fi
\noindent\textit{Source reference.} Original lemma 0F3R: sources/stacks-more-morphisms/pinned-original-source.tex, source lines 22966--22966. Navigation pointer only: the counted proof cites this statement by \textbackslash{}ref as a standing input; the deduction the counted proof has to make is carried out inside the counted span, and no part of that deduction is delegated to this pointer.\par\medskip
\ifdefined\hypertarget\hypertarget{sourcepointer-2}{}\fi
\noindent\textit{Source reference.} Original lemma 0F3L: sources/stacks-more-morphisms/pinned-original-source.tex, source lines 22749--22749. Navigation pointer only: the counted proof cites this statement by \textbackslash{}ref as a standing input; the deduction the counted proof has to make is carried out inside the counted span, and no part of that deduction is delegated to this pointer.\par\medskip
\ifdefined\hypertarget\hypertarget{sourcepointer-3}{}\fi
\noindent\textit{Source reference.} Original lemma 0F3K: sources/stacks-more-morphisms/pinned-original-source.tex, source lines 22731--22731. Navigation pointer only: the counted proof cites this statement by \textbackslash{}ref as a standing input; the deduction the counted proof has to make is carried out inside the counted span, and no part of that deduction is delegated to this pointer.\par\medskip
\ifdefined\hypertarget\hypertarget{sourcepointer-4}{}\fi
\noindent\textit{Source reference.} Original lemma 0F3B: sources/stacks-more-morphisms/pinned-original-source.tex, source lines 22173--22173. Navigation pointer only: the counted proof cites this statement by \textbackslash{}ref as a standing input; the deduction the counted proof has to make is carried out inside the counted span, and no part of that deduction is delegated to this pointer.\par\medskip
\ifdefined\hypertarget\hypertarget{sourcepointer-5}{}\fi
\noindent\textit{Source reference.} Original reference 01ZT: focused-seeds/geometry/S9-morph-asn-weighting/ref-tags.txt, source lines 5--5. Navigation pointer only: the counted proof cites this statement by \textbackslash{}ref as a standing input; the deduction the counted proof has to make is carried out inside the counted span, and no part of that deduction is delegated to this pointer.\par\medskip
\ifdefined\hypertarget\hypertarget{sourcepointer-6}{}\fi
\noindent\textit{Source reference.} Original reference 01SG: focused-seeds/geometry/S9-morph-asn-weighting/ref-tags.txt, source lines 6--6. Navigation pointer only: the counted proof cites this statement by \textbackslash{}ref as a standing input; the deduction the counted proof has to make is carried out inside the counted span, and no part of that deduction is delegated to this pointer.\par\medskip
\begin{proposition}
\label{proposition-asn-weighting}
Let $f : X \to Y$ be a surjective quasi-finite morphism of schemes.
Let $w : X \to \mathbf{Z}_{> 0}$ be a positive weighting of $f$.
Assume $X$ affine and $Y$ separated\footnote{It suffices if the
diagonal of $Y$ is affine.} and nonempty. Then the affine stratification
number of $Y$ is at most the number of distinct values of $\int_f w$
minus $1$.
\end{proposition}

\begin{proof}
Note that since $Y$ is separated, the morphism $X \to Y$ is affine
(Morphisms, Lemma \ref{morphisms-lemma-affine-permanence}).
The function $\int_f w$ attains its maximum $d$ by
Lemma \ref{lemma-max-int-w}. We will use induction on $d$.
Consider the open subscheme $Y_d = \{y \in Y \mid (\int_f w)(y) = d\}$
of $Y$ and recall that $f^{-1}(Y_d) \to Y_d$ is finite, see
Lemma \ref{lemma-max-int-finite}.
By Lemma \ref{lemma-affineness-of-large-open}
for every affine open $W \subset Y$ we have that $Y_d \cap W$ is affine
(this uses that $W \times_Y X$ is affine, being affine over $X$).
Hence $Y_d \to Y$ is an affine morphism of schemes. We
conclude that $f^{-1}(Y_d) = Y_d \times_Y X$ is
an affine scheme being affine over $X$.
Then $f^{-1}(Y_d) \to Y_d$ is surjective and
hence $Y_d$ is affine by Limits, Lemma \ref{limits-lemma-affine}.
If $Y_d = Y$, then the affine stratification number of $Y$ is $0$
and we conclude the lemma is true. If not, then let
$Y' = Y \setminus Y_d$ viewed as a reduced closed subscheme of $Y$.
Set $X' = Y' \times_Y X$.
Since $X'$ is closed in $X$ it is affine. Since
$Y'$ is closed in $Y$ it is separated.
The morphism $f' : X' \to Y'$ is surjective and
$w$ induces a weighting $w'$ of $f'$ (see
Lemma \ref{lemma-weighting-base-change}) and
$\int_{f'} w'$ is the restriction of $\int_f w$ to $Y'$.
By induction $Y'$ has an affine stratification of
length $\leq$ the number of distinct values of
$\int_{f'} w'$ minus $1$ and the proof is complete.
\end{proof}

\end{document}
